Combinatorial algorithms are often processed in large quantities, facilitated by file based data storage. 
An algorithm processes a class of objects from an input file and writes some output to yet another file.
Orbiter facilitates this kind of batch processing by means of input streams. 
The objects in one stream must all be of the same type.  
Table~\ref{tab:combinatorial:objects} shows the commands that are available to define input stream objects. 
\orbitertableofcommands{data_input_stream}
In Table~\ref{tab:combi:activity}, activities for combinatorial objects are listed.
\orbitertableofcommands{combinatorial_object_activity}
Some commands require extra options. 
Table~\ref{tab:combi:classification} shows options for classification.
\orbitertableofcommands{classification_of_objects}
Table~\ref{tab:combi:reporting} shows options for the report which is generated after a classification.
\orbitertableofcommands{classification_of_objects_report_options}

\bigskip

Let us consider some examples. 
First, we create the Hirschfeld surface.
Since the Hirschfeld surface is a cubic surface, 
the object is defined using point ranks in the relevant projective space 
as described in Section~\ref{sec:indexing:points}. 
For instance, the Hirschfeld surface in $\PG(3,4)$ has 45 points, 
and hence can be defined using a set of 45 codes of projective points.
The sequence of codes is recorded in a makefile variable:

\sourcecode{HIRSCHFELD_SURFACE_Q4_SET_OF_POINTS}


\sourcecode{Hirschfeld_q4_from_set}

The next example creates the two hyperovals in $\PG(2,16).$ 
The hyperovals are stored in makefile variables:

\sourcecode{HYPEROVAL_16_144}


\sourcecode{HYPEROVAL_16_16320}


\sourcecode{hyperoval_16_create}

In the next example, we read the points 
of an elliptic curve from file 
(see Section~\ref{sec:indexing:points}):

\sourcecode{EC_read}

In the next example, we read a packing, 
using a previously defined table of spreads, 
stored in a csv file.

\sourcecode{PG_3_5_assume_31_read}

The following command reads a file of large sets of designs:

\sourcecode{LS_AG_2_3_read}

The next command reads incidence geometries from a file containing the flags:

\sourcecode{geo_7_3_read}

The next command creates incidence geometries from a file containing row-ranks:

\sourcecode{Desargues_path_lex_least_read}


Combinatorial objects can be stored in text files. 
There can be any number of objects in one file.
The objects themselves are coded. 
For instance, a set of points in projective space 
is given as a set of integers, using the Orbiter point ranks. 
Likewise, a set of lines is coded using Orbiter line ranks. 
For designs, there are several ways in which the object can be stored. 
One way is by listing the incidences (flags) in a numerical form. 
An incidence between point $i$ and line $j$ is recorded as $ib+j.$ Here, $i$ and $j$ are the zero based index values of 
the point and line pair, and $b$ is the number of lines.
Another way is by recording the flags in a row of the incidence matrix as a subset. 
In this case, the number of incidences per row must be constant, say $r$. 
The flags in a given row thus comprise an $r$-subset of a set of size $b$, 
and the ranking of subsets is used to encode the flags in one row.
Proceeding in a row-by-row fashion, the incidence matrix is encoded as a vector of ranks of length $v$.
Yet another way is by listing the columns of the incidence matrix, provided the number of inciences is constant across columns, say $k$. 
In this case, we encode the flags in one column as a rank of a $k$-subset of a set of size $v$. 


\bigskip


The Pasch configuration is the configuration of 6 points and 4 lines 
shown below:
\orbiteroutput{GRAPHICS/WRAPPERS/pasch}
In the geometry, we have 4 lines, which we can identify with the index sets of the points as
$\{0,1,2\},$ $\{0,3,4\},$ $\{1,3,5\}$ and $\{2,4,5\}.$ 
The incidence matrix of the configuration is 
\orbiteroutput{GRAPHICS/LATEX/pasch_incma}
The rows correspond to the points, the columns correspond to the lines. 
The labels of points and lines are shown. 

\bigskip


How to encode an incidence structure for computer storage?
There are three ways to encode the incidence structure. 
One way encodes the flags of the geometry.
This will be described next.
The flag space is the set of all possible flags in the incidence matrix between the given number of points and lines.
The space is totally ordered 
using the row-major ordering
\orbiteroutput{GRAPHICS/LATEX/pasch_incma_row_major}
The Pasch configuration can now be coded as 
$$
\{0,1,4,6,8,11,13,14,17,19,22,23\}.
$$
The file \verb'pasch.inc' contains:
\orbiteroutput{GRAPHICS/LATEX/file_pasch_inc}
The first line lists the number of rows and columns of the incidence matrix, and the number of incidences. 
The geometry is encoded on the next line. After that, a marker of -1 shows that this is the only 
geometry in this file (the file format allows for any number of incidence geometries, all with the same parameters).
The final row is the order of the automorphism group of the geometry. 
This row is optional. In case that there are several geometries in the file, 
the orders will all be listed. In this case, the possible values will be 
collected with multiplicities, and listed in decreasing order.
The command 

\sourcecode{geo_pasch_read}

reads the incidence geometry from the file \verb'pasch.inc'.
It is also possible to enter the incidence geometry directly from the command line. The following example 
uses the \verb'-incidence_geometry' command to do so:

\sourcecode{geo_pasch_given}




